\chapter{Axionlike-particle-induced spin precession}\label{chap:theory}
\begin{jointwork}
    Part of the content in this chapter related to frequency-scanning considerations is adapted from Ref.\,\cite{Yuzhe2023_FrequencyScanning}.
\end{jointwork}

The ALP field can be detected through its gradient coupling to nuclear spins, which acts as an effective oscillating pseudomagnetic field. Understanding the properties of this coupling---its frequency, amplitude, and coherence---is essential for designing a sensitive experiment and choosing an optimal measurement strategy. This chapter builds the signal model from first principles: starting from the quantum-mechanical description of a single ALP and showing how the superposition of many such quanta gives rise to a classical stochastic oscillating field, which then drives resonant spin precession in the NMR sample.


\section{The ALP dark-matter field}\label{sec:theory:ALPfield}

\subsection{Single-particle wavefunction}\label{sec:theory:singleparticle}

In the non-relativistic limit appropriate for virialized dark matter ($v_a \ll c$), a single ALP with velocity $\mathbf{v}_a$ is described by a plane-wave state
\begin{equation}
    \psi(\mathbf{r}, t)
    = a_1\, e^{\,i\left(\mathbf{k}\cdot\mathbf{r} - 2\pi\nu_a t\right)}
    \,,\label{eq:single_ALP_wave}
\end{equation}
where the Compton frequency $\nu_a = m_a c^2/h$ is set by the rest mass alone, and the wavevector $\mathbf{k} = m_a \mathbf{v}_a/\hbar$ carries the kinetic information. The de~Broglie wavelength of such a particle is
\begin{equation}
    \lambda_{dB}
    = \frac{h}{m_a v_a}
    = \frac{hc}{m_a c^2}\cdot\frac{c}{v_a}
    \,.\label{eq:lambda_dB}
\end{equation}
For $m_a = \SI{1}{\micro\eV}$ and $v_a = \SI{220}{\km\per\s}$, this gives $\lambda_{dB} \approx \SI{1.7}{\km}$. The wavelength grows as $m_a^{-1}$, so at the neV masses probed in this thesis it reaches hundreds of kilometres---far larger than any laboratory apparatus. This means the ALP field is spatially coherent across the entire experiment, and the spatial gradient $\boldsymbol{\nabla}a$ that couples to nuclear spins is well approximated as uniform over the sample volume.

The enormous de~Broglie wavelength also implies an exceptionally high phase-space occupation number. The number density of ALPs in the local dark-matter halo is
\begin{equation}
    n_a = \frac{\rho_{DM}}{m_a c^2}
    \approx \frac{\SI{0.4}{\GeV\per\cm^3}}{m_a c^2}
    \,,\label{eq:na}
\end{equation}
so the mean occupation number per de~Broglie mode is
\begin{equation}
    \mathcal{N} \approx n_a\,\lambda_{dB}^3
    = \frac{\rho_{DM}}{m_a c^2}\left(\frac{hc}{m_a c^2}\right)^3\!\left(\frac{c}{v_a}\right)^3
    \,.\label{eq:occupation}
\end{equation}
For $m_a = \SI{1}{\micro\eV}$ and $\rho_{DM} = \SI{0.4}{\GeV\per\cm^3}$, one obtains $\mathcal{N} \approx 10^{30}$. Even for the lightest masses still in the particle regime ($m_a \sim \SI{e-22}{\eV}$, ``fuzzy'' dark matter), $\mathcal{N}$ is astronomically large. This is a direct consequence of the tiny ALP mass: for a fixed energy density, a lighter boson must be exponentially more abundant.

\subsection{Classical-field limit from mode superposition}\label{sec:theory:classicallimit}

Because ALPs are bosons and their occupation numbers are so large, quantum statistics strongly suppress number fluctuations relative to the mean: $\delta\mathcal{N}/\mathcal{N} \sim \mathcal{N}^{-1/2} \to 0$. In this regime the quantum field operator $\hat{a}$ can be replaced by a classical $c$-number field---the same approximation that turns a quantum electromagnetic field into a classical wave when the photon number is large.

The dark-matter halo does not consist of a single momentum mode but of a superposition drawn from the Maxwell-Boltzmann velocity distribution. Each mode $i$ contributes a plane wave with amplitude $a_i$, wavevector $\mathbf{k}_i = m_a \mathbf{v}_i/\hbar$, and a random phase $\phi_i$:
\begin{equation}
    a(\mathbf{r}, t)
    = \sum_i a_i\, \cos\!\bigl(\nu_i t - \mathbf{k}_i\cdot\mathbf{r} + \phi_i\bigr)
    \,,\label{eq:ALP_sum}
\end{equation}
where $\nu_i = \nu_a + \tfrac{1}{2}m_a v_i^2/h \approx \nu_a$ is the total energy (including kinetic energy). On time scales short compared to the spread of phases accumulating from the slightly different $\nu_i$, the sum behaves as a single coherent oscillation at frequency $\nu_a$:
\begin{equation}
    a(\mathbf{r},t) = a_0 \cos(\nu_a t - \mathbf{k}\cdot \mathbf{r} + \phi)
    \,.\label{eq:ALP_field}
\end{equation}
Here $\mathbf{k} = m_a \mathbf{v}_a/\hbar$ is an effective wavevector, $\phi$ is a random phase uniformly distributed in $[0,\,2\pi)$, and the amplitude $a_0$ follows a Rayleigh distribution because it is the magnitude of a sum of many random phasors \cite{Centers2021Stochastic}. Its r.m.s.\ value is set by the local dark-matter energy density \cite{BERGSTROM1998137,Jungman1992SupersymmetricDM,doi:10.1146/annurev.astro.39.1.137}:
\begin{align}
    \rho_{DM} \approx \dfrac{c\, m_a^2 a_0^2}{2\hbar^3}
     \,.\label{eq:rhoDM}
\end{align}

\subsection{Coherence time}\label{sec:theory:stochasticity}

The modes in Equation~\eqref{eq:ALP_sum} have slightly different frequencies due to the kinetic-energy spread. Their phases drift apart on a time scale set by the spectral width of the field. For a Maxwell-Boltzmann velocity distribution with average speed $v_a \approx \SI{220}{\km\per\s}$ \cite{Evans2019_SHMRefined}, the ALP field spectral linewidth (FWHM) observed by a detector on Earth is
\begin{equation}
    \Gamma_a = \dfrac{\nu_a}{Q_a}
    \,,\label{eq:Gamma_a_Qa}
\end{equation}
where $Q_a \equiv (c/v_a)^2 \approx 1.8\times10^6$ is the ALP quality factor. The ALP coherence time $\tau_a$, defined by $\Gamma_a = 2\pi/\tau_a$, is therefore
\begin{equation}
    \tau_a = \dfrac{2\pi Q_a}{\nu_a}
    \,.\label{eq:taua}
\end{equation}
Within one coherence time the superposition Equation~\eqref{eq:ALP_sum} is well approximated by the single coherent mode Equation~\eqref{eq:ALP_field}. On longer time scales the random phases between modes cause the amplitude and phase of the effective oscillation to drift, so the field must be treated as a stochastic process \cite{Centers2021Stochastic}. For Compton frequencies above \SI{1}{\kHz}, $\tau_a \lesssim \SI{1000}{\second}$; the interplay between $\tau_a$ and the NMR relaxation times determines the optimal measurement strategy in Chapter~\ref{chap:results}.


\section{Periodic modulations}

\subsection{Daily modulation}

The ALP-field gradient driving spin precession depends on the direction of the relative velocity between the Earth and the local dark-matter halo. As the Earth rotates, this direction sweeps through space with a period of one sidereal day. An ALP signal therefore exhibits a daily modulation of its amplitude and phase \cite{Gramolin2022Feb_SpectralSignatures}.

\subsection{Annual modulation}

In addition to daily rotation, the Earth's orbital velocity around the Sun modulates the relative ALP-detector velocity with a period of one year. Both the daily and annual modulations have precisely predictable periods and phases, making them useful consistency checks for any claimed dark-matter-candidate signal \cite{Gramolin2022Feb_SpectralSignatures}.


\section{Spin-precession signal}\label{sec:theory:signal}

The nuclear-gradient coupling of the ALP field to nuclear spins is described by the Hamiltonian
\begin{align}
    H_{\mathrm{aNN}} \sim g_{\mathrm{aNN}}\, {\boldsymbol \nabla} a \cdot \mathbf{I}
   \,,\label{eq:HaNN}
\end{align}
where $g_{\mathrm{aNN}}$ is the coupling strength in units of $\SI{}{\GeV^{-1}}$, ${\boldsymbol \nabla} a$ is the gradient of the ALP field and $\mathbf{I}$ is the nuclear spin operator. This coupling acts as a pseudomagnetic drive field for nuclear spins precessing at the Larmor frequency in a bias field $\mathbf{B}_0$. Its r.m.s.\ Rabi frequency is \cite{Centers2021Stochastic,Yuzhe2023_FrequencyScanning}
\begin{equation}
    \Omega_a \approx \dfrac{1}{2}\, g_{\mathrm{aNN}} \sqrt{2\hbar c\, \rho_{DM}}\, v_{\bot}
    \,,\label{eq:Omega_a}
\end{equation}
where $v_{\bot}$ is the component of the ALP-detector relative velocity perpendicular to $\mathbf{B}_0$. Taking $v_{\bot} = \SI{220}{\km\per\second}$, we have $\Omega_a/2\pi = g_{\mathrm{aNN}} \times \SI{0.2}{\GeV\cdot\Hz}$. For $g_{\mathrm{aNN}} \lesssim 10^{-10}\,\SI{}{\GeV^{-1}}$, the Rabi period $2\pi/\Omega_a \gtrsim 1000\,\mathrm{yr}$, confirming that the ALP field is a weak drive on the spins.


\subsection{Relaxations and coherence}\label{sec:theory:relaxations}

The transverse spin relaxation in the sample arises from homogeneous mechanisms (e.g., spin-spin interactions), characterised by the time $T_2$, and from inhomogeneous mechanisms (e.g., magnetic-field gradients across the sample), which cause additional dephasing. The total effective transverse relaxation time $T_2^* \leq T_2$ incorporates both contributions.

The NMR linewidth (FWHM, assuming a Lorentzian lineshape) of the free-induction-decay spectrum is
\begin{equation}
    \Gamma_{n} = \dfrac{2}{T_2^*}
    \,,\label{eq:Gamma_n}
\end{equation}
and the normalized linewidth is
\begin{equation}
    \delta_\nu \equiv \dfrac{\Gamma_{n}}{\nu} = \dfrac{2}{\nu T_2^*}
    \,,\label{eq:delta_nu}
\end{equation}
often expressed in parts per million (\SI{}{\ppm}). The inhomogeneous contribution to $1/T_2^*$ can be reduced by magnetic shimming, giving control over $T_2^*$ independently of the intrinsic $T_2$.


\subsection{Estimation}\label{sec:theory:estimation}

For a sample with equilibrium longitudinal magnetization $M_0$ prepared along $\mathbf{B}_0$, the ALP-field drive tilts the magnetization by a small angle $\xi \ll 1$ (weak-drive regime, $\Omega_a T \ll 1$), generating a transverse magnetization $M_1 \approx M_0\xi$. The r.m.s.\ transverse magnetization is \cite{Yuzhe2023_FrequencyScanning}
\begin{equation}
    \sqrt{\langle M_{1}^2 \rangle} \approx \dfrac{1}{\sqrt{2}}\, u_n M_{0} \sqrt{\langle \xi^2 \rangle}
    \,,\label{eq:M1_rms}
\end{equation}
where $u_n$ is the fraction of nuclear spins on resonance with the ALP field. Approximating both the NMR and ALP lineshapes as rectangles,
\begin{equation}
    u_n \approx \left\{
    \begin{array}{cl}
    1\,,
    &  \tau_a \ll T_2^* \\[6pt]
    \dfrac{\Gamma_a}{\Gamma_{n}} = \dfrac{\pi T_2^*}{\tau_a}\,,
    &  T_2^* \ll \tau_a
    \end{array} \right.
    \,.\label{eq:un}
\end{equation}

The r.m.s.\ tipping angle $\sqrt{\langle\xi^2\rangle}$ is limited by the interplay between ALP decoherence and spin relaxation. Modeling the ALP decoherence as a random phase that resets on the time scale $\tau_a$, the tipping angle executes a two-dimensional random walk in the rotating frame with step size $\Omega_a\tau_a$ \cite{Yuzhe2023_FrequencyScanning}. Combining this with the three relevant relaxation scales,
\begin{equation}
    \sqrt{\langle \xi^2\rangle}  = \left\{
    \begin{array}{cl}
    \Omega_{a}\sqrt{\tau_a T_2^*} ,
    &  \tau_a \ll T_2^*\\[6pt]
    \Omega_a\dfrac{\tau_a}{\sqrt{\pi}} ,
    &  T_2^*\ll\tau_a \ll T_2\\[6pt]
    \Omega_a T_2 ,
    &  T_2 \ll \tau_a
    \end{array} \right.
    \,.\label{eq:rms_xi}
\end{equation}
These cases combine into the r.m.s.\ transverse magnetization (for dwell time $T \gg \tau_a,\,T_2^*,\,T_2$):
\begin{equation}
    \sqrt{\langle M_{1}^2 \rangle} \approx \dfrac{M_{0}\Omega_a}{\sqrt{2}}\left\{
    \begin{array}{cl}
    \sqrt{\tau_a T_2^*}\, ,
    &  \tau_a \ll T_2^*\\[6pt]
    \dfrac{T_2^*}{\sqrt{\pi}}\, ,
    &  T_2^*\ll\tau_a \ll T_2\\[6pt]
    \dfrac{T_2^* T_2}{\sqrt{\pi}\,\tau_a}\, ,
    &  T_2 \ll \tau_a
    \end{array} \right.
    \,.\label{eq:M1_regimes}
\end{equation}
In the third case, when $T_2 \ll \tau_a$ and the dwell time $T \gg \tau_a$, the signal amplitude decreases because the ALP field cannot drive the spins coherently for longer than $T_2$, while off-resonance spins (those not within the narrow ALP linewidth) contribute negligibly.


\subsection{Numerical simulation}\label{sec:theory:simulation}

The analytical estimates above are valid in the limit $T \gg \tau_a,\,T_2^*,\,T_2$. At lower Compton frequencies, the ALP coherence time can become comparable to or longer than experimentally accessible dwell times, and the simple scaling laws no longer apply. Numerical simulations using the \texttt{axionbloch} package (Appendix~\ref{app:axionbloch}) are used to validate the analytical estimates and to explore intermediate parameter regimes.

Three representative situations arise:
\begin{description}
    \item[Low frequency] $\tau_a \gg T$: the ALP field is essentially coherent over the measurement. The signal grows linearly in time and the $T^{-1/4}$ sensitivity scaling is not yet reached.
    \item[Medium frequency] $\tau_a \approx T$: the transition regime, treated numerically.
    \item[High frequency] $T \gg \tau_a,\,T_2^*,\,T_2$: the analytical results apply. For hyperpolarized \ch{^{129}Xe}, $T_1$ relaxation limits the achievable polarization, introducing an additional reduction factor of approximately $1/3$ relative to the fully polarized case.
\end{description}


\subsection{Summary}\label{sec:theory:summary}

Depending on the relative magnitudes of $T_2$, $T_2^*$, and $\tau_a$, four regimes are identified \cite{Yuzhe2023_FrequencyScanning}:
\begin{enumerate}[label = (\arabic*)]
    \item $\tau_a \ll T_2^*$;
    \item $T_2^* \ll \tau_a \ll T_2 $;
    \item $T_2^* \ll T_2 \ll \tau_a$;
    \item $T_2^*=T_2 \ll \tau_a$.
\end{enumerate}
These regimes are illustrated in Figure~\ref{fig:T2star_T2_taua_regimes}. The key parameters for each regime are summarised in Table~\ref{tab:summary_of_regimes}, where $\nu_{\mathrm{start}}$ and $\nu_{\mathrm{end}}$ are the lower and upper limits of the target frequency band.

\begin{figure}[htb]
    \centering
    \includegraphics[width=.6\linewidth]{works/frequency-scanning/T2star_T2_taua_relations_20230413.png}
    \caption{Regimes of the spin-precession experiment based on the relative magnitudes of $T_2$, $T_2^*$, and $\tau_a$. Reproduced from Ref.\,\cite{Yuzhe2023_FrequencyScanning}.}
    \label{fig:T2star_T2_taua_regimes}
\end{figure}

\begin{table}[ht]
    \centering
    \caption{Summary of parameters for the four regimes. The normalised linewidth $\delta_\nu = \tau_a/(\pi Q_a T_2^*)$ is assumed constant in each regime. The optimal dwell time is derived in Appendix~\ref{app:dwell_time}.}
    \begin{tabular}{cccccc}
        \hline
        & Regime & Optimal dwell time $T$ & Signal linewidth $\Gamma$ & Spectral factor $u_n$ & r.m.s.\ tipping angle $\sqrt{\langle \xi^2\rangle}$ \\
        \hline
        & & & & & \\
        (1) & $\tau_a \ll T_2^*$
            & $\dfrac{T_{\mathrm{tot}}}{Q_a \ln\!\left( \nu_{\mathrm{end}}/\nu_{\mathrm{start}} \right)}$
            & $\dfrac{2}{T_2^*}$ & $1$ & $\Omega_{a}\sqrt{\tau_a T_2^*}$\\
        & & & & & \\
        (2) & $T_2^* \ll \tau_a \ll T_2$
            & $\dfrac{T_{\mathrm{tot}}\,\tau_a}{\pi Q_a T_2^*\ln\!\left( \nu_{\mathrm{end}}/\nu_{\mathrm{start}} \right)}$
            & $\dfrac{2\pi}{\tau_a}$ & $\dfrac{\pi T_2^*}{\tau_a}$ & $\Omega_a\dfrac{\tau_a}{\sqrt{\pi}}$ \\
        & & & & & \\
        (3) & $T_2^* \ll T_2 \ll \tau_a$
            & $\dfrac{T_{\mathrm{tot}}\,\tau_a}{\pi Q_a T_2^*\ln\!\left( \nu_{\mathrm{end}}/\nu_{\mathrm{start}} \right)}$
            & $\dfrac{2\pi}{\tau_a}$ & $\dfrac{\pi T_2^*}{\tau_a}$ & $\Omega_a T_2$\\
        & & & & & \\
        (4) & $T_2^*=T_2 \ll \tau_a$
            & $\dfrac{T_{\mathrm{tot}}\,\tau_a}{\pi Q_a T_2\ln\!\left(\nu_{\mathrm{end}}/\nu_{\mathrm{start}}\right)}$
            & $\dfrac{2\pi}{\tau_a}$ & $\dfrac{\pi T_2}{\tau_a}$ & $\Omega_a T_2$ \\
        & & & & & \\
        \hline
    \end{tabular}
    \label{tab:summary_of_regimes}
\end{table}
